\documentclass[12pt,a4paper]{article}
\usepackage[margin=2.5cm]{geometry}
\usepackage{xcolor}
\usepackage{amsmath,amssymb}
\usepackage{graphicx}
\usepackage{hyperref}
\usepackage{booktabs}

% Paleta de colores de NaTex
\definecolor{foxorange}{RGB}{255, 107, 53}

\hypersetup{
    colorlinks=true,
    linkcolor=foxorange,
    urlcolor=foxorange
}

\title{
    \vspace{-1cm}
    \includegraphics[width=4.5cm]{logo_natex.png}\\[0.5cm]
    {\Huge \bfseries \color{foxorange} Guia Rapida de NaTex Studio}\\[0.2cm]
    {\large \textbf{Version 0.1 - Editor Cientifico y Colaborativo}}
}
\author{\textbf{NaTex Team} \\ Tu nuevo entorno de LaTeX sin complicaciones}
\date{\today}

\begin{document}

\maketitle

\begin{abstract}
\textbf{Bienvenido a NaTex Studio v0.1!} Este documento sirve tanto de bienvenida como de instructivo practico interactivo. NaTex fue disenado para estudiantes, profesores e investigadores que desean la potencia total de LaTeX combinada con la inmediatez de la previsualizacion local, dictado de clases por voz, colaboracion en tiempo real y asistencia por Inteligencia Artificial.
\end{abstract}

\section{Atajos de Teclado Esenciales}
El flujo de trabajo en NaTex esta optimizado para la velocidad. Memoriza estos atajos clave:

\begin{table}[h]
\centering
\begin{tabular}{@{}llp{8.5cm}@{}}
\toprule
\textbf{Atajo} & \textbf{Accion} & \textbf{Descripcion} \\ \midrule
\texttt{Ctrl + S} & \textbf{Compilar} & Genera y recarga instantaneamente el visor PDF a la derecha. \\
\texttt{Alt + M} & \textbf{Formulas} & Abre el Constructor y Teclado Matematico Cientifico. \\
\texttt{Tab} & \textbf{Indentar} & Inserta 4 espacios en el cursor para ordenar el codigo. \\
\texttt{Ctrl + Enter} & \textbf{Compilar Rapido} & Alternativa rapida para compilar sin despegar las manos del teclado. \\
\bottomrule
\end{tabular}
\end{table}

\section{Constructor de Ecuaciones (Alt + M)}
Olvidate de buscar comandos raros de memoria. Presiona \textbf{Alt + M} (o clic en el boton \textit{Formulas}) para usar el panel visual interactivo con vista previa en vivo KaTeX.

Ejemplo de formula cuantica insertada con el asistente:
\begin{equation}
    i\hbar \frac{\partial}{\partial t}|\psi(t)\rangle = \hat{\mathbf{H}}|\psi(t)\rangle
\end{equation}

Ejemplo de calculo multivariable:
\begin{equation}
    \oint_{C} \vec{B} \cdot d\vec{r} = \mu_0 I_{\text{enc}} + \mu_0 \varepsilon_0 \frac{\partial \Phi_E}{\partial t}
\end{equation}

\section{Dictafono de Clases para Movil}
Estas en clase o en una reunion? 
\begin{enumerate}
    \item Haz clic en el boton \textbf{"Clase"} en la barra superior (o en Inicio).
    \item Escanea el codigo QR con tu celular (no requiere instalar ninguna aplicacion).
    \item Presiona \textbf{Grabar} mientras el profesor explica.
    \item NaTex procesara el audio con \textbf{Google Gemini} y creara un apunte LaTeX estructurado con formulas matematicas, teoremas y resumen ejecutivo listo para compilar.
\end{enumerate}

\section{Galeria y Gestion de Figuras}
Para insertar imagenes en tus documentos:
\begin{itemize}
    \item Arrastra cualquier archivo (.png, .jpg, .svg) a la barra lateral izquierda o usa el menu \textbf{Insertar $\rightarrow$ Galeria de Imagenes}.
    \item Haz clic en el boton \textbf{Copiar LaTeX} para copiar el snippet \texttt{\textbackslash begin\{figure\}} y pegalo donde quieras.
    \item Puedes renombrar archivos directamente y NaTex actualizara tu codigo de forma automatica.
\end{itemize}

\begin{figure}[h]
    \centering
    \includegraphics[width=0.32\textwidth]{logo_natex.png}
    \caption{Logotipo oficial de NaTex Studio.}
    \label{fig:logo}
\end{figure}

\section{Colaboracion en Tiempo Real}
Haz clic en el menu \textbf{Colaborar $\rightarrow$ Compartir con un amigo}. NaTex generara un tunel seguro con un enlace y codigo QR. Tu companero podra entrar desde cualquier lugar del mundo para editar el documento y compilar en vivo contigo sin configuraciones de red complejas.

\section{Empieza a crear!}
Puedes borrar este texto o modificarlo como quieras para empezar tu propio documento. Disfruta redactando con NaTex!

\end{document}